\chapter{系统错误、内核崩溃与可观测性}

系统会出错、会 panic、会变慢。本章把错误路径收敛、崩溃恢复和可观测性工具放在一起：从 errno 到 kdump，从 printk 到 eBPF，从现象到证据到结论。

\section{错误、oops 与 panic}

操作系统必须区分普通错误、进程错误、内核 bug 和硬件不可恢复错误。普通错误应返回 \code{errno}，例如文件不存在、权限不足、非阻塞 I/O 暂时不可用；进程错误应向进程发送信号，例如非法地址访问；内核 bug 可能触发 WARN、oops 或 panic；硬件错误可能要求隔离设备、下线页面或重启系统。

\begin{warningbox}
若把所有错误都 panic，系统不可用；若把所有错误都吞掉，数据会被悄悄破坏。错误分类的责任在于明确每个失败点：哪些返回 \code{errno}，哪些终止进程，哪些停止系统。
\end{warningbox}

panic 的含义是内核认为继续运行比停止更危险。典型原因包括核心不变量破坏、调度器状态自相矛盾、文件系统元数据检测到不可恢复损坏、内核栈溢出、无法处理的机器检查。panic 不是调试打印的高级版本，而是系统级 fail-stop 策略：停止服务，尽量保留证据，避免进一步破坏状态。

\begin{keyidea}
高质量 OS 实现不只写正常路径，还要把每个失败点的责任写清楚：返回错误、终止进程、重试、降级、隔离，还是 panic。
\end{keyidea}

错误路径的工程难点在资源回收。系统调用可能已经分配页、获取锁、增加引用、插入表项、发起 I/O，然后第七步失败。正确实现必须让失败路径按照相反顺序撤销已完成步骤，并且不能撤销还没完成的步骤。C 语言内核常用 \code{goto out\_...} 标签表达这种线性回滚；C++ 内核或教学模拟器可以用 RAII 表达资源所有权，但必须注意中断上下文、panic 路径和 no-throw 约束。

转向可观测性：操作系统问题很少只靠读代码解决。线程为什么不运行，可能是调度等待、锁等待、I/O 阻塞、page fault、信号、cgroup 限流或下游服务；内存为什么上涨，可能是真实泄漏、allocator 缓存、page cache、mmap、线程栈或内核 buffer；写文件为什么慢，可能是 page cache 脏页、设备队列、fsync、目录同步或日志提交。没有观测证据，结论只是猜测。

可观测性要把 OS 状态映射到证据。进程和线程状态可以从 \filepath{/proc}、ps、top、调度 trace 观察；系统调用可以用 strace 观察；CPU 热点可以用采样 profiler；page fault、context switch、I/O 等待和网络状态可以用系统计数器和 tracing 工具；应用层必须补充业务 trace，记录请求身份、队列等待、状态机阶段和错误码。

\begin{keyidea}
可观测性的目标不是收集越多越好，而是把 OS 状态映射到证据：on-CPU 解释运行时消耗，off-CPU 解释等待时间，两者闭合才能定位问题。
\end{keyidea}

调试时要区分 on-CPU 和 off-CPU。on-CPU 说明线程正在执行指令；off-CPU 说明线程没有运行，可能在等待。CPU profiler 只能解释运行时消耗在哪里，解释不了大部分时间都在睡眠的请求。尾延迟问题常常要从 off-CPU 等待开始看。

内存调试要区分 virtual size、resident set、anonymous memory、file-backed page、page cache、allocator arena 和 kernel memory。\code{malloc} 后不触页，可能只增加虚拟地址；首次写入才分配物理页；释放内存后 RSS 不马上下降，也可能是 allocator 保留 arena。

\begin{warningbox}
把所有 RSS 增长都叫泄漏，是不合格诊断。虚拟地址、物理页、page cache、allocator arena 和内核内存各有不同来源，必须先分清再下结论。
\end{warningbox}

I/O 调试要区分 buffered I/O、direct I/O、mmap、page cache hit/miss、dirty writeback、fsync 和设备吞吐。一次写很快，可能只是进入 page cache；一次读很快，可能只是热缓存；一次 fsync 很慢，可能是在偿还之前累计的脏页。实验必须控制冷缓存和热缓存，否则数字不可比较。

\section{从症状建立证据链}

错误分类可以写成下表：

\begin{longtable}{|L{0.18\linewidth}|L{0.36\linewidth}|L{0.34\linewidth}|}
\toprule
类别 & 处理方式 & 示例 \\
\midrule
可预期用户错误 & 返回 \code{errno} & \code{ENOENT}、\code{EACCES}、\code{EAGAIN} \\\tablerowrule
用户进程违规 & 信号或终止该进程 & page fault 权限错误、非法指令 \\\tablerowrule
资源暂时不足 & 返回错误、等待或回收 & \code{ENOMEM}、队列满、端口耗尽 \\\tablerowrule
设备局部故障 & 重置设备、隔离队列、上报 I/O error & 磁盘超时、网卡 reset \\\tablerowrule
内核不变量破坏 & WARN/oops/panic & 双重释放、runqueue 计数错误 \\\tablerowrule
持久化一致性风险 & remount read-only、replay journal、panic & 元数据校验失败 \\
\bottomrule
\end{longtable}

一个系统调用实现应在设计阶段列出失败点：

\begin{lstlisting}
open(path, flags):
  copied_path = copy_path_from_user(path)        may fail EFAULT/ENAMETOOLONG
  parent = resolve_parent(copied_path)           may fail ENOENT/EACCES
  inode = lookup_or_create(parent, name, flags)  may fail ENOENT/EEXIST/EROFS
  file = alloc_file(inode, flags)                may fail ENOMEM
  fd = alloc_fd(current.files)                   may fail EMFILE
  install_fd(fd, file)                           commit point
  return fd
\end{lstlisting}

提交点之前失败，要释放已拿到的引用；提交点之后失败，通常不能再返回“没有打开”，而要让 fd 表保持一致或执行明确关闭。把提交点写出来，可以避免半初始化 fd 留在表里。

本章的实践模板是“现象、假设、证据、反证、结论”。

\begin{longtable}{|L{0.18\linewidth}|L{0.72\linewidth}|}
\toprule
字段 & 请求延迟升高示例 \\
\midrule
现象 & p95 从 20ms 升到 500ms \\\tablerowrule
假设 & worker 等锁或等 I/O \\\tablerowrule
证据 & trace 中 queue\_wait 低，blocked\_io 高，线程睡在 writeback \\\tablerowrule
反证 & on-CPU profiler 没有热函数占用 500ms \\\tablerowrule
结论 & 延迟主要来自 I/O 完成等待，不是 CPU 算法变慢 \\
\bottomrule
\end{longtable}

程序卡死时，先列线程状态：哪个线程持有锁，哪个线程等待条件变量，哪个线程在 join，关闭标志是否设置，是否有线程仍可能唤醒等待者。若所有消费者都睡在 not\_empty，而关闭路径只通知 not\_full，就能定位到等待谓词和通知对象不匹配。

实践中要保留原始证据：命令、时间、输入规模、机器信息、内核版本、容器限制、采样频率、日志片段和无法观察的字段。没有这些，报告不可复现。

\section{指标、事件、调用链与转储}

\subsection{errno 返回与用户可见契约}

\code{errno} 不是随便选的数字。调用者会根据错误码决定重试、降级、提示用户还是终止。

\begin{lstlisting}
copy_from_user_checked(dst, user_ptr, len):
  if not user_range_valid(user_ptr, len, READ):
    return EFAULT
  if len > MAX_COPY:
    return EINVAL
  fault = copy_may_fault(dst, user_ptr, len)
  if fault:
    return EFAULT
  return OK
\end{lstlisting}

把非法用户地址返回 \code{EFAULT}，而不是 panic，是因为这是用户进程可触发的预期错误。内核必须把“不可信输入”当正常情况处理。只有当内核自己的页表、锁状态或对象引用被破坏时，才进入更严重路径。

\subsection{goto 回滚模板}

C 风格内核常见模式：

\begin{lstlisting}
sys_create(path):
  name = null
  inode = null
  file = null

  name = copy_name(path)
  if name == null:
    ret = ENOMEM
    goto out

  inode = create_inode(name)
  if is_error(inode):
    ret = error(inode)
    goto out_name

  file = alloc_file(inode)
  if file == null:
    ret = ENOMEM
    goto out_inode

  ret = install_fd(file)
  if ret < 0:
    goto out_file
  return ret

out_file:
  put_file(file)
out_inode:
  put_inode(inode)
out_name:
  free(name)
out:
  return ret
\end{lstlisting}

这个结构的优点是回滚顺序与获取顺序相反，且每个标签只释放已经成功获取的资源。缺点是标签命名和跳转边界容易腐烂，需要测试每个失败注入点。C++ RAII 可以减少样板，但内核中不是所有路径都适合异常；析构函数也不能在持有自旋锁或 panic 路径中做复杂阻塞操作。

\subsection{WARN、BUG、oops 与 panic}

可以把严重程度分成四级：

\begin{longtable}{|L{0.18\linewidth}|L{0.32\linewidth}|L{0.36\linewidth}|}
\toprule
机制 & 含义 & 处理 \\
\midrule
WARN & 发现异常但可能继续 & 打印栈、计数、继续或降级 \\\tablerowrule
BUG/oops & 当前内核路径不可继续 & 杀死当前进程或停止当前 CPU，视上下文而定 \\\tablerowrule
panic & 整个系统不可安全继续 & 停机、重启或进入 crash dump \\\tablerowrule
watchdog & 系统长时间无响应 & 触发诊断、中断栈采样或 panic \\
\bottomrule
\end{longtable}

教学系统可以实现 \code{assert\_kernel\_invariant}：调试构建 panic，发布构建记录错误并尽量隔离。但对文件系统元数据损坏这类风险，继续写入可能扩大破坏，remount read-only 或 panic 是合理选择。

\subsection{看门狗与故障转储}

watchdog 检测“没有前进”。硬锁死可能表现为中断不再发生；软锁死可能表现为某 CPU 长时间不调度；RCU stall 表示读侧临界区过久导致回收无法完成。

\begin{lstlisting}
watchdog_tick(cpu):
  now = time()
  if now - cpu.last_schedule_time > SOFTLOCKUP_LIMIT:
    dump_cpu_stack(cpu)
    warn("soft lockup")
  if now - cpu.last_interrupt_time > HARDLOCKUP_LIMIT:
    panic("hard lockup")
\end{lstlisting}

panic 后的目标不是恢复服务，而是保留证据：panic 原因、CPU 寄存器、栈、当前任务、锁持有、最近日志、脏页状态和设备错误。crash dump 的价值在于让下一次启动后能分析，而不是在 panic 路径中做复杂修复。

\subsection{重试与幂等}

不是所有失败都应该立即返回。临时内存不足可以触发回收后重试；设备超时可以 reset 后重试；网络发送可以等待窗口。但重试必须有边界，并且操作要么幂等，要么有明确的提交记录。

\begin{lstlisting}
retry_io(req):
  while req.retries < MAX_RETRIES:
    submit(req)
    status = wait_completion(req, timeout)
    if status == OK:
      return OK
    if not is_transient(status):
      return status
    reset_path_if_needed(req.device)
    req.retries++
  return EIO
\end{lstlisting}

对写请求，若设备状态不明，简单重试可能造成重复写。对普通块覆盖写通常可接受；对“追加一条记录”这类上层操作，必须有序列号、校验或日志来识别重复提交。

\subsection{事件日志}

最简单的可观测性算法是事件日志。每个关键状态转换追加一条结构化记录：时间、线程、请求 ID、事件类型、对象 ID、旧状态、新状态和错误码。日志不能只写自然语言，否则无法聚合和验证。

\begin{lstlisting}
record_event(req_id, object, old_state, new_state, reason):
  entry.ts = monotonic_time()
  entry.cpu = current_cpu()
  entry.thread = current_thread()
  entry.req_id = req_id
  entry.object = object
  entry.old_state = old_state
  entry.new_state = new_state
  entry.reason = reason
  ring_append(trace_buffer, entry)
\end{lstlisting}

事件日志本身也要有背压。无限日志会把故障变成磁盘或内存问题。常见做法是固定大小环形缓冲区、按级别采样、按请求 ID 打开详细追踪，或者把关键错误同步写入可靠日志。

\subsection{延迟分解}

请求延迟不能只记录总耗时。一个可行动的分解至少包括 queue wait、runnable wait、on CPU、blocked on lock、blocked on I/O、downstream wait 和 cleanup。实现上可以在状态转换处打点，然后按请求 ID 聚合。

\begin{lstlisting}
transition(req, new_state):
  now = clock()
  req.time_in_state[req.state] += now - req.last_ts
  req.state = new_state
  req.last_ts = now
\end{lstlisting}

这个算法的复杂度是每次状态转换 \code{O(1)}，但要求状态枚举稳定。若所有等待都记成 blocked，后续无法判断是锁、I/O 还是 page fault。

\subsection{off-CPU 分析}

很多性能问题不在 CPU 火焰图里，因为线程大部分时间不在 CPU 上。off-CPU 分析记录线程从运行态切到睡眠态的原因，以及被谁唤醒。

\begin{lstlisting}
on_schedule_out(task, reason, wait_object):
  task.offcpu_start = clock()
  task.wait_reason = reason
  task.wait_object = wait_object

on_wakeup(task, waker):
  delta = clock() - task.offcpu_start
  offcpu_histogram[task.wait_reason].add(delta)
  record_wakeup_edge(waker, task, task.wait_object)
\end{lstlisting}

这个数据能区分等锁、等 I/O、等页、等网络、等定时器和等子进程。若 p99 延迟主要来自 off-CPU 的 writeback 等待，优化业务 CPU 代码不会解决问题。

\subsection{资源对账}

资源泄漏调试需要对账算法。每次 open 增加 fd 计数，每次 close 减少；每次分配对象记录 owner，每次释放删除记录；系统退出或测试结束时检查剩余对象。

\begin{lstlisting}
track_alloc(id, type, owner):
  live[id] = {type, owner, stack}

track_free(id):
  live.erase(id)

assert_no_leak():
  for each object in live:
    report(object)
\end{lstlisting}

真实系统不能给所有对象都保存完整 stack，因为成本太高；可以采样、按类型开启、只在测试构建开启。但思想不变：资源必须能被计数，生命周期必须能闭合。

\subsection{最小可复现报告}

高质量调试报告要让别人能复现判断。最小格式如下：

\begin{lstlisting}
Problem:
  现象、影响范围、首次出现时间

Environment:
  内核版本、硬件/容器限制、文件系统、负载规模

Timeline:
  关键事件按时间排序

Evidence:
  系统调用、trace、线程状态、资源计数、日志、指标

Rejected hypotheses:
  已经排除的解释和证据

Conclusion:
  当前最能解释现象的机制

Boundary:
  仍未验证的假设和下一步实验
\end{lstlisting}

报告的价值不在于一次猜中，而在于缩小不确定性。写出反证能防止团队反复讨论已经排除的方向。

\section{故障注入与因果验证}

可观测性与错误路径都需要测试。一个系统若只能在成功时输出日志，失败时没有状态，就不可维护。

\begin{enumerate}
\tightlist
\item errno 矩阵：非法指针、权限不足、资源耗尽返回不同错误码。
\item 失败注入：在系统调用每个分配点失败，引用计数和锁计数最终归零。
\item 提交点测试：提交前失败无用户可见副作用，提交后状态自洽。
\item WARN 测试：非致命不变量异常记录证据但不破坏后续测试。
\item panic 测试：核心不变量破坏后停止调度，输出原因和栈。
\item watchdog 测试：构造关中断死循环或不调度循环，能触发诊断。
\item 设备重试：临时错误可重试，永久错误返回 \code{EIO} 且唤醒等待者。
\item crash dump：panic 后保存当前任务、寄存器、锁和最近日志。
\end{enumerate}

\begin{itemize}
\tightlist
\item 每个长时间等待点都能输出等待原因或 trace 事件。
\item 每个请求都有稳定身份，能串起进入队列、开始执行、阻塞、完成和失败。
\item 错误码保留原始原因，不被统一改写成 unknown error。
\item 资源计数可对账：打开 fd 数、活跃线程数、队列长度、in-flight I/O、分配对象数。
\item 性能测试区分冷启动、热缓存、预热、输入规模和测量次数。
\item 故障注入能触发日志和恢复路径，而不是只测试成功路径。
\end{itemize}

本章要求读者能够区分 errno、信号、oops 与 panic，为系统调用写出失败回滚路径，并解释
提交点、幂等重试和故障证据的关系。诊断结论还必须能够复现：若报告只写“怀疑是系统问题”，
却没有线程状态、系统调用、等待点、资源计数或反证，就不能构成操作系统层面的分析。

\section{调度、内存、I/O 与网络诊断}

\subsection{BUG、WARN、oops 和 panic 的代码路径}

Linux 中 \code{WARN\_ON} 记录调用栈但通常继续；\code{BUG\_ON} 表示当前路径不可继续；
oops 是内核异常报告，可能终止当前进程或进一步导致 panic；\code{panic} 停止系统或按配置
重启。报告需要收集寄存器、调用栈、污染标志、已加载模块和当前任务，才能为离线分析保留上下文。

\begin{lstlisting}
bad_kernel_access:
  fault in kernel mode
  -> do_error_trap or page fault path
  -> oops_begin()
  -> show registers, stack, code bytes
  -> decide die, panic, or kill current task
\end{lstlisting}

错误严重程度取决于上下文。用户进程传坏指针应返回 \code{EFAULT}；内核在持有自旋锁时 page fault，可能说明严重 bug；文件系统检测到元数据不一致继续写入可能扩大损坏，因此 remount read-only 或 panic 都是合理策略。

\subsection{KASAN、KCSAN 和 KFENCE}

\begin{longtable}{|L{0.18\linewidth}|L{0.44\linewidth}|L{0.26\linewidth}|}
\toprule
工具 & 发现问题 & 取舍 \\
\midrule
KASAN & 越界、use-after-free、部分栈/全局错误 & 高开销，高覆盖 \\\tablerowrule
KCSAN & 数据竞争，非原子并发访问 & 采样式，可能需多跑 \\\tablerowrule
KFENCE & 堆越界和 UAF & 低开销，覆盖有限 \\
\bottomrule
\end{longtable}

这些工具不会替代锁设计。它们帮助你发现违反不变量的具体证据，例如哪个对象释放后被哪个路径继续访问，哪个字段被两个 CPU 无同步写入。测试报告应保留工具输出中的对象地址、分配栈、释放栈和访问栈。

\subsection{fault injection}

Linux 提供多种故障注入，例如 fail\_page\_alloc、failslab、fail\_make\_request。目标是在可控位置让分配、slab、块 I/O 失败，验证错误路径。

\begin{lstlisting}
test plan:
  fail next kmalloc in sys_open
  expect no fd installed, inode ref dropped

  fail block request number N during fsync
  expect fsync returns EIO and writeback error recorded
\end{lstlisting}

没有故障注入，错误回滚代码很可能从未执行过。高质量 OS 代码要让每个 \code{goto out\_*} 标签都能被测试触发。故障注入测试必须检查“没有发生什么”：没有泄漏引用、没有遗留锁、没有半初始化对象进入全局表。

\subsection{kdump 与 /proc/vmcore}

kdump 预留一段崩溃捕获内存。panic 后通过 kexec 启动捕获内核，导出内存转储，再结合
符号信息分析。panic 路径越少依赖已经损坏的系统状态，转储成功率越高。

\begin{lstlisting}
panic
  -> stop other CPUs
  -> save registers and notes
  -> kexec crash kernel
  -> expose /proc/vmcore
  -> analyze task, locks, stacks, memory
\end{lstlisting}

生产环境中，panic 策略要结合服务目标：有的系统宁可快速重启，有的系统必须保留 dump，有的存储系统检测到元数据损坏要停止写入避免扩大破坏。

\subsection{错误恢复的状态机写法}

恢复不是“再试一次”。每个可恢复对象都应有明确状态：正常、降级、恢复中、不可用、已隔离。状态转换要有进入条件、退出条件和超时策略。

\begin{lstlisting}
device state:
  ONLINE
    on timeout -> RECOVERING
  RECOVERING
    stop queues
    reset hardware
    replay configuration
    complete lost requests
    if success -> ONLINE
    if fail count exceeded -> FAILED
  FAILED
    reject new requests with EIO
    keep diagnostic evidence
\end{lstlisting}

没有状态机，恢复路径容易和正常 I/O 并发踩踏：reset 时仍提交新请求，失败请求没有唤醒等待者，或恢复成功后旧 completion 又回来释放同一资源。

\subsection{错误预算和降级}

并非所有错误都要立即 panic 或无限重试。系统可以设置错误预算：短时间少量校验失败触发重试，连续失败触发隔离，核心元数据错误触发只读或 panic。错误预算让系统避免因为瞬时故障过度反应，也避免因为反复失败无限拖延。

\begin{longtable}{|L{0.22\linewidth}|L{0.34\linewidth}|L{0.32\linewidth}|}
\toprule
现象 & 初始处理 & 超过预算后 \\
\midrule
块 I/O 超时 & 重试、reset 队列 & 下线设备或返回 \code{EIO} \\\tablerowrule
内存分配失败 & 回收、压缩、等待 & OOM kill 或返回 \code{ENOMEM} \\\tablerowrule
网络重传增多 & 降低发送速率 & 熔断连接或限流 \\\tablerowrule
文件系统校验失败 & 停止相关写入 & remount read-only 或 panic \\
\bottomrule
\end{longtable}

错误预算必须可观测：重试次数、最后错误、进入降级时间、当前状态和影响对象都要暴露。否则系统处在降级状态时，用户只会看到“偶尔很慢”。

\subsection{cleanup attribute 与 RAII 的边界}

用户态 C 可以用 \code{__attribute__((cleanup))} 模拟作用域清理，C++ 可以用 RAII。内核 C 常用 \code{goto out}，是因为错误路径需要精确控制上下文、锁和释放顺序。三种方法没有绝对高低，关键是所有权必须清楚。

\begin{lstlisting}
// conceptual RAII style
FileRef file = open_checked(path)
PageRef page = allocate_page()
LockGuard guard(inode.lock)
commit_update(file, page)
\end{lstlisting}

RAII 适合普通进程上下文中的资源组合；但析构函数不应隐藏可能睡眠的复杂 I/O，也不应在 panic 路径中试图获取新锁。教学 OS 可以使用 C++20 RAII 管理引用和锁，但要为中断上下文、noexcept 和显式提交点写清规则。

\subsection{panic 报告模板}

panic 输出要稳定、短而有证据。最低字段如下：

\begin{lstlisting}
panic report:
  reason
  cpu id and current task
  instruction pointer and stack pointer
  trap vector or syscall number
  held locks
  taint or loaded modules
  last N kernel log records
  affected device or inode if any
  reboot/dump policy
\end{lstlisting}

报告不是越长越好。panic 时系统可能已经损坏，输出过多可能再次阻塞。核心原则是保留能定位第一现场的事实，而不是在崩溃路径中做复杂分析。

\subsection{printk 和 dmesg}

\code{printk} 有 loglevel、rate limiting、console 输出和 ring buffer。panic 前后的日志可能是唯一证据，但过量 printk 会扰动时序，甚至拖慢系统。驱动错误路径应记录设备、请求 id、状态寄存器、errno，而不是只写“failed”。

\subsection{ftrace 和 tracepoint}

ftrace 可以追踪函数入口、函数图和静态 tracepoint。tracepoint 的价值是稳定字段：调度切换、irq、block、net、syscalls 都有可解析事件。

\begin{lstlisting}
trace-cmd record -e sched:sched_switch -e block:block_rq_issue
trace-cmd report
\end{lstlisting}

trace event 需要定义稳定字段，写入每 CPU 环形缓冲区，再由用户态读取。tracepoint 的
作用是把一次状态变化转化为可采集证据；字段定义和丢失策略决定这些证据能回答哪些问题。

\subsection{perf 与 PMU}

\code{perf stat} 看计数，\code{perf record/report} 看采样调用栈，\code{perf top} 看实时热点。PMU 事件能观察 cycles、instructions、cache misses、branch misses、TLB misses。IPC 低可能来自 cache miss、分支错误、前端供给不足或后端等待，不能只说“CPU 慢”。

\begin{lstlisting}
perf stat -e cycles,instructions,cache-misses,branches,branch-misses ./workload
perf record -g ./workload
perf report
\end{lstlisting}

\subsection{eBPF 和 bpftrace}

eBPF 程序可挂到 kprobe、uprobe、tracepoint、profile、LSM、cgroup、网络路径。BPF map 保存状态，helper 提供受控内核访问。bpftrace 适合快速一行观测。

\begin{lstlisting}
bpftrace -e 'tracepoint:syscalls:sys_enter_write { @[comm] = count(); }'
bpftrace -e 'kprobe:vfs_read { @[kstack] = count(); }'
\end{lstlisting}

eBPF 的边界是 verifier、安全权限和观测开销。生产系统应限制采样频率和输出量，避免观测本身制造故障。

\subsection{strace、ltrace 和 gdb}

strace 基于 ptrace 观察系统调用 enter/exit，能回答“程序在请求内核做什么”；ltrace 观察动态库调用，能回答“用户态库在做什么”；gdb/kgdb 用符号和断点观察状态。strace 看不到内核内部等待原因，perf/ftrace/eBPF 可以补上。

\subsection{/proc、/sys、sysctl 和 systemd journal}

\filepath{/proc/meminfo}、\filepath{/proc/interrupts}、\filepath{/proc/slabinfo}、\filepath{/proc/[pid]/status}、\filepath{/proc/[pid]/stack} 是运行状态窗口；\filepath{/sys} 暴露 kobject、设备、cgroup、block queue 参数；\filepath{/proc/sys} 是 sysctl 控制面；journald 给系统日志结构化字段。

\subsection{PSI：Pressure Stall Information}

PSI 记录 CPU、memory、IO stall 时间，回答“任务因为资源压力停顿了多久”。它比单纯利用率更接近用户感受到的延迟。内存 PSI 高说明任务因回收或内存不足停顿；I/O PSI 高说明任务等待 I/O；CPU PSI 高说明 runnable 但等不到 CPU。

\begin{lstlisting}
cat /proc/pressure/cpu
cat /proc/pressure/memory
cat /proc/pressure/io
\end{lstlisting}

PSI 适合做过载保护信号。服务可以在 memory/io pressure 高时限流，避免把队列堆到超时。

\subsection{工具选择决策表}

可观测性工具很多，选择时先问问题类型。

\begin{longtable}{|L{0.26\linewidth}|L{0.42\linewidth}|L{0.2\linewidth}|}
\toprule
问题 & 优先工具 & 证据 \\
\midrule
程序卡在哪个系统调用 & \code{strace -tt -T} & syscall、耗时、errno \\\tablerowrule
CPU 时间花在哪里 & \code{perf record -g} & on-CPU 调用栈 \\\tablerowrule
线程为什么不运行 & off-CPU trace、sched trace & 睡眠原因、唤醒者 \\\tablerowrule
磁盘请求慢在哪里 & block trace、iostat、perf sched & 队列和 completion \\\tablerowrule
内存压力是否存在 & PSI、\filepath{/proc/meminfo}、vmstat & stall、reclaim、swap \\\tablerowrule
内核路径偶发错误 & tracepoint、kprobe、bpftrace & 事件字段和栈 \\\tablerowrule
容器被限制 & cgroup v2 文件、journal & throttled、OOM、pids limit \\
\bottomrule
\end{longtable}

\begin{warningbox}
不要一上来抓所有数据。先用低成本指标定位层次，再用高精度 trace 验证假设。观测也会扰动系统，尤其是高频 kprobe、过量 printk 和全量系统调用 trace——把被观察系统变成了另一个系统。
\end{warningbox}

\subsection{一次延迟诊断的完整演示}

假设 HTTP 服务 p99 从 80ms 升到 2s。诊断不应直接改代码，而应建立证据链。

\begin{lstlisting}
step 1: application trace
  parse: 1ms, business: 5ms, persist: 1900ms

step 2: strace
  fsync(fd) = 0 <1.87s>

step 3: /proc/pressure/io
  some avg10 increased

step 4: block trace
  many flush requests wait behind writeback

step 5: conclusion
  tail latency dominated by storage flush and dirty writeback
\end{lstlisting}

反证也要写：CPU profiler 没有 2s on-CPU 热点；网络 recv/send 没有长时间阻塞；锁等待直方图没有对应峰值。这样结论才不是“看起来像磁盘”。

\subsection{指标设计：counter、gauge、histogram、event}

系统指标至少分四类。

\begin{longtable}{|L{0.2\linewidth}|L{0.34\linewidth}|L{0.34\linewidth}|}
\toprule
类型 & 含义 & OS 示例 \\
\midrule
counter & 单调增加事件数 & syscall 次数、错误次数、重试次数 \\\tablerowrule
gauge & 当前状态值 & 队列长度、活跃线程、脏页数量 \\\tablerowrule
histogram & 分布 & 请求延迟、锁等待、I/O 完成时间 \\\tablerowrule
event & 结构化状态转换 & 任务睡眠、唤醒、请求提交、panic \\
\bottomrule
\end{longtable}

错误率只用 counter 不够，还要有错误码和对象维度；延迟只报平均值不够，要有 p50/p95/p99 和最大值；队列只报当前值不够，最好有高水位和丢弃次数。

\subsection{低开销 instrumentation 原则}

观测点应贴近状态转换，而不是在循环里无节制打印。推荐做法：

\begin{itemize}
\tightlist
\item 热路径用 per-CPU counter，避免全局锁。
\item 高频事件采样或按条件开启。
\item 错误路径记录结构化字段和稳定 id。
\item 日志消息带 request id、task id、object id。
\item trace buffer 有固定容量和丢弃计数。
\item 观测代码本身不能持有业务关键锁太久。
\end{itemize}

可观测性也是代码，必须有性能预算。一个 tracing patch 若把锁竞争扩大，就可能把被观察系统变成另一个系统。

\subsection{验收：可复现诊断包}

本卷要求最终报告能够打包复现。最小诊断包包括：

\begin{lstlisting}
diagnostic bundle:
  exact command and workload
  kernel version and hardware/container limits
  application logs with request ids
  relevant /proc and /sys snapshots
  trace or perf data
  benchmark raw results
  analysis notes with rejected hypotheses
\end{lstlisting}

报告不是把命令输出堆在一起，而是把证据组织成因果链：现象是什么，在哪层出现，哪些解释被排除，当前结论的边界在哪里。能做到这一点，才算真正掌握 OS 调试。
